\section{Continuation to shrinking spheres}\label{sec:balls}

We now apply the product continuation criterion near a point where an
isometry is already known on one side.  The purpose of this section is to
construct a family of functionals on harmonic functions, represented on
small spheres.  Their existence will follow from the geometry of the
spheres.  A uniform bound for the functionals, needed when the spheres
shrink, will follow from the estimate in the next section.

\subsection{The local data and their normalization}
\label{balls:setup}

Throughout Sections~\ref{sec:balls}--\ref{sec:bootstrap}, fix the
manifolds $(N_j,g_j)$, actual Dirichlet Green functions $G_j$, common
harmonic coordinate neighborhood $\Omega$, matching set $W$, and finite
harmonic pairs $(u_1^a,u_2^a)$ of Theorem~\ref{local:extension}.
Thus the coordinate metrics and paired functions agree on $W$, and the
Green kernels agree on $W\times W$ off the diagonal. The first-side
coordinate Hessians span the hyperplane annihilated by $g_1^{-1}(0)$. An interior ball in
$W$ tangent at $0$ supplies the one-sided hypersurface required by
that theorem. We construct the sphere functionals for these data.

A linear change of harmonic coordinates makes $g_1(0)=g_2(0)=I$ and
writes the separating hypersurface as a graph tangent to $x_n=0$, with
the known side below it.  For a spatial
dilation $\lambda>0$, write
\begin{equation}\label{balls:dilation}
 g_{j,\lambda}(x)=g_j(\lambda x),\qquad
 G_{j,\lambda}(x,y)=\lambda^{n-2}G_j(\lambda x,\lambda y).
\end{equation}
Dilate the harmonic pairs at the same time:
\[
 u_{j,\lambda}^a(x)=u_j^a(\lambda x).
\]
They remain harmonic for $g_{j,\lambda}$ and agree on the dilated known
set. The metrics and Green kernels in \eqref{balls:dilation} have the
corresponding normalization on the dilated neighborhoods. On each fixed
bounded coordinate set the metrics converge smoothly to $I$ as $\lambda\downarrow0$.  The dilated known set
contains every fixed compact set lying strictly below $x_n=0$, once
$\lambda$ is sufficiently small.

Use $x$ for the first variable and $y$ for the second.  The mixed kernel
$\mathcal G$ is initially $G_{1,\lambda}$ when $y$ is in the known set,
and $G_{2,\lambda}$ when $x$ is in that set.  These prescriptions agree on
their overlap and solve the first equation in $x$ and the second in $y$,
away from the diagonal.  Proposition~\ref{cross:initial} continues this
kernel to any specified compact region of fixed positive separation,
with smooth bounds uniform for small $\lambda$.  We shall use that result
only after all the positive separation ranges in question have been fixed.

The geometry of a moving sphere dictates the next choices. In the
strict extension test, the leading terms involving the gradient of its
radius cancel; strict concavity of the logarithmic radius supplies the
remaining positive term (see \eqref{balls:test-expression}). We therefore
choose a conformal reference metric for which a depth function is
strictly concave. Separately, a scalar normalization of the second
inverse metric will give the two inverse matrices the same trace
relative to the reference metric. The harmonic Hessians will determine
their remaining difference.

Fix the coordinate cylinder
\[
 Y=\{y=(y',y_n): |y'|\leq1,\ -1\leq y_n\leq1\},
 \qquad t_0(y)=y_n+|y'|^2.
\]
Take a sufficiently large fixed number $L$ and put
\begin{equation}\label{balls:normalization}
 \gamma=e^{-2Lt_0}g_{1,\lambda},\qquad
 Q=\frac{n\,g_{2,\lambda}^{-1}}
          {\operatorname{tr}(\gamma g_{2,\lambda}^{-1})},\qquad
 h=Q-\gamma^{-1}.
\end{equation}
Thus $Q$ is the principal matrix of a positive multiple of the second
equation and $\operatorname{tr}(\gamma h)=0$.  The first equation can
likewise be multiplied by a positive function to give principal metric
$\gamma$.  Such multiplication does not change its solutions.  The
metric $\gamma$ will also identify the two tangent spaces at a common
coordinate point.  All these coefficients have uniform smooth bounds
on a fixed neighborhood of $Y$ for sufficiently small $\lambda$, and
$h\to0$ in every fixed smooth norm as $\lambda\downarrow0$.

Define the depth function
\begin{equation}\label{balls:depth}
 t(y)=\int_{1/8}^{t_0(y)}e^{-2Ls}\,ds.
\end{equation}
It is negative at $0$ and has nonvanishing differential on $Y$.  The
conformal connection formula gives
\begin{equation}\label{balls:depth-hessian}
 \operatorname{Hess}_{\gamma}t
 =e^{-2Lt_0}\bigl(
    \operatorname{Hess}_{g_{1,\lambda}}t_0
       -L|dt_0|_{g_{1,\lambda}}^2g_{1,\lambda}\bigr)
 \leq-c\gamma.
\end{equation}
Indeed $|dt_0|$ is bounded below, whereas the first Hessian on the right is
uniformly bounded.  Choose $L$ first and then an upper bound for
$\lambda$ so that the displayed strict inequality holds with fixed $c>0$.

\subsection{Radii and cutoff regions}

For a fixed integer $p\geq2$, a positive amplitude $R_*$, and
$0<\delta\leq1$, define
\begin{equation}\label{balls:radius}
 \ell_\delta(t)=\delta\log(1+e^{t/\delta}),\qquad
 r_\delta(y)=R_*\ell_\delta(t(y))^p,
 \qquad f_\delta=\log r_\delta,\quad
 \beta_\delta=|df_\delta|_\gamma.
\end{equation}
The balls will be the $\gamma$-geodesic balls with center $y$ and radius
$r_\delta(y)$, lying in the first coordinate neighborhood. As the lemma
below shows, they collapse where $t\le0$, including a neighborhood of
the contact point, but their radii stay bounded below where $t$ is
bounded away from zero on the positive side. This distinction will keep
cutoff derivatives away from uncontrolled shrinking spheres.

\begin{lemma}[Profile bounds]\label{balls:profiles}
On $Y$, after the preceding choices of geometry, the following bounds
hold uniformly for small $\lambda$, $0<\delta\leq1$, and $p\geq2$:
\begin{equation}\label{balls:profile-estimates}
 \beta_\delta\geq c p,\qquad
 \operatorname{Hess}_\gamma f_\delta\leq-c\beta_\delta\gamma,
 \qquad |\partial^j f_\delta|\leq C_j\beta_\delta^j
 \quad(j\geq1).
\end{equation}
The constants $c,C_j$ do not depend on $p$.  The functions
$r_\delta^{1/p}$ have a Lipschitz bound independent of $\delta$ and
$\lambda$ for fixed $p,R_*$.  The radii increase with $\delta$, tend to
zero on $\{t\leq0\}$ as $\delta\downarrow0$, and satisfy
\begin{equation}\label{balls:radius-smallness}
 \sup_{\delta,y}r_\delta\beta_\delta^2\longrightarrow0
 \quad\text{as }R_*\downarrow0
\end{equation}
for each fixed $p$.
\end{lemma}

\begin{proof}
Write $a_\delta=(\log\ell_\delta)'$.  In terms of $s=t/\delta$,
\[
 a_\delta(t)=\delta^{-1}A(s),\qquad
 A(s)=\frac{e^s}{(1+e^s)\log(1+e^s)}.
\]
The function $A$ is positive and decreasing.  At the negative end it
tends to one, and at the positive end it is comparable to $(1+s)^{-1}$.
Differentiating the formula, or its convergent expressions at the two
ends, gives $|A^{(j)}(s)|\leq C_jA(s)^{j+1}$.  On the remaining compact
interval the same inequalities follow by positivity.  On the fixed
$t$ interval this gives a positive lower bound for $a_\delta$ and
\[
 a_\delta'\leq0,\qquad
 |\partial_t^j\log\ell_\delta|\leq C_j a_\delta^j,
 \qquad a_\delta\leq\ell_\delta^{-1}.
\]
Since $|dt|_\gamma$ is bounded above and below,
$\beta_\delta=p a_\delta|dt|_\gamma$.  Equation
\eqref{balls:depth-hessian} and
\[
 \operatorname{Hess}_\gamma f_\delta
   =p a_\delta\operatorname{Hess}_\gamma t
          +p a_\delta'\,dt\otimes dt
\]
prove the Hessian bound.  The chain rule proves the other derivative
bounds, with constants independent of $p\geq2$ because
$p a_\delta^j\leq p^j a_\delta^j$.

The bound $0<\ell_\delta'<1$ gives the asserted Lipschitz estimate for
$r_\delta^{1/p}=R_*^{1/p}\ell_\delta\circ t$.  Moreover,
\[
 \partial_\delta\ell_\delta
  =\log(1+e^s)-\frac{s e^s}{1+e^s}>0.
\]
The function on the right increases up to $s=0$ and decreases afterward,
with limit zero at both ends.  The pointwise limit of $\ell_\delta$ is
$\max(t,0)$.  Finally,
\[
 r_\delta\beta_\delta^2
       \leq C R_*p^2\ell_\delta^{p-2},
\]
and $\ell_\delta$ is uniformly bounded on the fixed interval.  This
proves \eqref{balls:radius-smallness}.
\end{proof}

Choose $t_b>0$ so small that $3t_b<t(1/4)$, where on the right $t$ denotes
the increasing function of $t_0$ in \eqref{balls:depth}, and set
\begin{equation}\label{balls:K}
 K=\{y\in Y:y_n\geq-1/2,\ t(y)\leq2t_b\}.
\end{equation}
Then $K\Subset Y^\circ$.  We fix a smooth cutoff $\chi\in
C_c^\infty(Y^\circ)$ equal to one on a neighborhood of $K$, with
\begin{equation}\label{balls:cutoff-region}
 \operatorname{supp}(d\chi)\subset
       \{y_n<-1/4\}\cup\{t>t_b\}.
\end{equation}
For example, multiply a cutoff that changes from zero to one between
$y_n=-3/4$ and $y_n=-1/2$ by a cutoff of $t$ that changes from one to zero
between $2t_b$ and $3t_b$, moving the endpoints slightly to make it one
on a neighborhood of $K$.  The inequality $3t_b<t(1/4)$ keeps its support
away from the lateral and upper faces of $Y$.  This gives fixed room
around $K$ and places every cutoff derivative either in the known lower
region or where
\[
 r_\delta\ge R_*t_b^p>0.
\]
Thus the later estimates for the sphere functionals can use the original
Green kernel or fixed-separation continuation on every cutoff derivative
support.

\subsection{The strict test at a moving sphere}

For coordinate derivatives of the tensor $h$, use the pointwise notation
\[
 H_8(y)=\sum_{|\alpha|\leq8}|\partial_y^\alpha h(y)|.
\]
We shall temporarily impose
\begin{equation}\label{balls:bootstrap-condition}
 H_8\leq\varepsilon r_\delta,\qquad
 r_{\delta'}\beta_{\delta'}^2\leq\varepsilon
       \quad(\delta\leq\delta'\leq1)
 \quad\text{on }Y.
\end{equation}
The first inequality measures how closely the two principal metrics
agree compared with the current radius.  It will be improved in
Section~\ref{sec:bootstrap}.

\begin{proposition}[Continuation to the sphere family]\label{balls:continuation}
Choose $p$ sufficiently large and then $\varepsilon>0$ sufficiently small
in terms of the fixed background bounds.  Choose $R_*>0$ small enough
that all the balls lie in normal neighborhoods and
\eqref{balls:radius-smallness} is at most $\varepsilon$.
For all sufficiently small spatial dilations $\lambda$, and every
$0<\delta\le1$ for which \eqref{balls:bootstrap-condition} holds, the mixed
kernel extends to a neighborhood of
\[
 \{(x,y):y\in Y,\ d_\gamma(x,y)=r_\delta(y)\}.
\]
The extension solves the two homogeneous equations and agrees with the
prescribed lower cross and the fixed-separation continuation.  It is
smooth for each positive $\delta$; no uniform norm as $\delta\downarrow0$
is asserted here.
\end{proposition}

\begin{proof}
We check Proposition~\ref{cross:criterion} at each intermediate radius
$r=r_{\delta'}$.  Put $f=\log r$, $\beta=|df|_\gamma$, and
$b=\nabla^\gamma r$, so $|b|=r\beta$.  Use the defining function
\[
 \rho(x,y)=\tfrac12\bigl(d_\gamma(x,y)^2-r(y)^2\bigr).
\]
The positive side is the outside of the sphere.  Parallel translation
along its radial geodesic identifies the two tangent spaces for
$\gamma$; write $\omega$ for the unit vector from $y$ to $x$.

In the normalization of the strict test, the normal and tangential
vectors can be combined into complex vectors
\[
 U_1=\omega+i v,\quad v\perp\omega,\quad |v|=1,
 \qquad U_2=q+i w,
\]
where the conditions for the second principal metric $Q^{-1}$ give
\begin{equation}\label{balls:test-vectors}
 (\omega+b)\cdot U_2=1,\qquad
 q=\frac{\omega+b}{|\omega+b|^2}+O(|h|),\qquad
 |w|=|\omega+b|^{-1}+O(|h|).
\end{equation}
Here and below lengths and dot products in these formulas use $\gamma$
and its parallel identification.  The first equation includes
$(\omega+b)\cdot w=0$.  Since $|b|$ and $|h|$ are small, both partial
gradients of $\rho$ are nonzero and $|U_2|$ is bounded above and below.

The Hessian of half the squared distance, with the product
$\gamma$-connection, evaluated on these endpoint vectors is
\[
 |U_1-U_2|^2+O(r^2)(|U_1|^2+|U_2|^2).
\]
To see the error size, parametrize the radial geodesic on $[0,1]$.
Its Jacobi field with prescribed endpoint values differs, in a parallel
frame, from the affine interpolant by $O(r^2)$ in $C^1$, by the Jacobi
equation and bounded curvature.  The second variation formula then gives
the displayed bound.  Replacing the second connection by that of
$Q^{-1}$ costs $O(r^2)$ as well: the connection difference is
$O(|h|+|\partial h|)=O(\varepsilon r)$ and $|d\rho|=O(r)$.

Now
$\omega\cdot(U_1-U_2)=b\cdot U_2$ and
$\operatorname{Hess}(r^2/2)=2\,dr\otimes dr+r^2\operatorname{Hess}f$.
Consequently the strict Hessian expression is
\begin{equation}\label{balls:test-expression}
 |\pi_{\omega^\perp}(U_1-U_2)|^2-|b\cdot U_2|^2
      -r^2\operatorname{Hess}_\gamma f(U_2,\overline U_2)+O(r^2).
\end{equation}
For a real Hessian, its value on a complex vector and its conjugate is
the sum of its values on the real and imaginary parts, as required by
the criterion.

The real part of the projected difference is
$-\pi_{\omega^\perp}b+O(|b|^2+|h|)$.  For its imaginary part,
\eqref{balls:test-vectors} gives
$|\pi_{\omega^\perp}w|=1-\omega\cdot b+O(|b|^2+|h|)$;
compare this with $|v|=1$.  Squaring these two bounds yields
\[
 |\pi_{\omega^\perp}(U_1-U_2)|^2
 \geq |b|^2-C\bigl(|b|^3+|b||h|+|h|^2\bigr).
\]
The real and imaginary directions of $U_2$ are almost orthonormal, so
\[
 |b\cdot U_2|^2\leq
       \bigl(1+C(|b|+|h|)\bigr)|b|^2.
\]
By Lemma~\ref{balls:profiles}, the expression
\eqref{balls:test-expression} is therefore bounded below by
\[
 c r^2\beta
  -C\bigl(r^2+r^3\beta^3+\varepsilon r^2\beta
                       +\varepsilon^2r^2\bigr).
\]
Choose $p$ so that $\inf\beta$ absorbs the $Cr^2$ term.  Then choose
$\varepsilon$ small.  The inequality $r\beta^2\leq\varepsilon$ absorbs
$r^3\beta^3$, and the expression is strictly positive.

It remains to specify the starting and boundary data for this local
test.  Deform $\delta'$ from $1$ to $\delta$, using $(y,\omega)$ as base
variables.  The inequality $H_8\leq\varepsilon r_\delta$ implies its
counterpart for every $\delta'\geq\delta$, because the radii increase
with $\delta'$.  At $\delta'=1$ the required separations have a fixed
positive minimum.  The same is true wherever $t>t_b$, since
$r_{\delta'}\geq R_*t_b^p$.  Proposition~\ref{cross:initial} supplies
uniformly bounded data on these fixed separation ranges, with open room.
Where $y_n<-1/4$, use the original $G_{1,\lambda}$; for small $\lambda$
this is part of the known lower cross.  These regions cover the lateral
and end boundaries of the cylinder that are not supplied by the initial
parameter, and their prescriptions agree on overlaps.

The base sphere bundle over the closed cylinder is compact.  On each
fixed interval $\delta\leq\delta'\leq1$, the radii are positive, all
surfaces lie in normal neighborhoods, and the strict tests just proved
apply away from the prescribed regions.  Lemma~\ref{cross:propagation}
therefore supplies compatible extension germs on the whole family.
Its finite-cover gluing also retains agreement with the prescribed
regions.  This is an existence argument for each finite interval; it
does not assert uniform continuation bounds when the lower endpoint
tends to zero.
\end{proof}

\subsection{Functionals represented on the spheres}

The mixed kernel now defines a map from first-metric harmonic functions
to second-metric harmonic functions.  Its representation on a small
sphere will allow us to estimate this map without estimating the kernel
through every continuation step.
Fix $0<\delta\le1$ and write $r=r_\delta$ in this subsection.

For a function $u$ harmonic for $g_{1,\lambda}$ near
$\overline{B_\gamma(y,r(y))}$, define
\begin{equation}\label{balls:flux}
 T_yu=\int_{\partial B_\gamma(y,r(y))}
  \bigl(\mathcal G(x,y)\partial_{\nu_1}u(x)
       -u(x)\partial_{\nu_1}\mathcal G(x,y)\bigr)\,dS_1(x).
\end{equation}
Here $\nu_1$ is the outer unit normal and $dS_1$ the hypersurface measure
for the actual first metric $g_{1,\lambda}$, not for $\gamma$.
Let $S_y$ be the unit sphere of $(T_y\mathbb R^n,\gamma_y)$, with its
rotation-invariant probability measure $d\sigma_y$.  Parametrize the
integration sphere by $x=\exp_y^\gamma(r(y)\omega)$.

\begin{proposition}[The transfer density]\label{balls:functional}
Under the assumptions of Proposition~\ref{balls:continuation}, there is
a smooth function $\psi(y,\omega)$ on the unit sphere bundle over
$Y^\circ$ such that
\begin{equation}\label{balls:density-representation}
 T_yu=\int_{S_y}\psi(y,\omega)
              u\bigl(\exp_y^\gamma(r(y)\omega)\bigr)\,d\sigma_y(\omega).
\end{equation}
For every fixed center $y_0$, if $u$ is first-metric harmonic on a
neighborhood of $\overline{B_\gamma(y_0,r(y_0))}$, then
$y\mapsto T_yu$ is defined and second-metric harmonic on a sufficiently
small neighborhood of $y_0$.  This neighborhood may depend on $u$.
In particular, for functions harmonic on the fixed coordinate
neighborhood required by the sphere family,
\begin{equation}\label{balls:reproduction}
 \int_{S_y}\psi\,d\sigma_y=1,
 \qquad T_yx^i=y^i,
 \qquad T_yu_{1,\lambda}^a=u_{2,\lambda}^a(y).
\end{equation}
On the support of derivatives of $\chi$, every fixed angular smooth
norm of $\psi$ and its first base derivatives is bounded independently
of $\delta$ and small $\lambda$, after $p,R_*$ and the fixed geometry
have been chosen.
\end{proposition}

\begin{proof}
For each ball, its smooth Dirichlet solution operator determines the
normal derivative of a harmonic function from its boundary trace.
Move this boundary operator in the first term of \eqref{balls:flux}
onto the smooth boundary value of $\mathcal G$ by adjunction, and then
pull back to $S_y$.  The resulting smooth density is $\psi$.
Smooth dependence of the Dirichlet problem and of the sphere
parametrization gives its smooth dependence on $y$ for each positive
radius.  This construction initially applies to smooth boundary
values and hence to every harmonic function used in
\eqref{balls:density-representation}.

Fix a center $y_0$ and a function $u$ harmonic on a neighborhood of its
closed ball.  The continued kernel is defined on an open neighborhood
of the compact sphere over $y_0$.  Choose a fixed slightly larger
surrounding surface inside this kernel neighborhood and the domain
of $u$.  After restricting the center neighborhood, it surrounds every
moving sphere and the whole intervening collar remains in that common
domain.  Green's identity on the collar replaces the moving integral
in \eqref{balls:flux} by the integral over this fixed surface, since
both first-variable equations are homogeneous there.  Applying the
second-variable operator under the fixed integral proves the stated
local harmonicity.  This argument uses arbitrary local harmonic
functions, not only the source-generated pairs.

On the lower known region $\mathcal G=G_{1,\lambda}$, so Green's formula
gives $T_yu=u(y)$.  Unique continuation on the connected set $Y^\circ$
now proves all identities in \eqref{balls:reproduction}, first for the
constant and coordinate pairs and then for each given harmonic pair.

For the uniform bounds on the cutoff region, distinguish the two sets
in \eqref{balls:cutoff-region}.  On the known lower set, $T_y$ is ordinary
evaluation.  Its density is the Poisson density at the center of the
ball, expressed in normalized coordinates.  After pulling the equation
back by $z\mapsto\exp_y^\gamma(rz)$ and multiplying by the Laplacian
scaling, the coefficients form a smooth uniformly elliptic family down
to $r=0$, whose limiting equation is the Euclidean Laplace equation.
The interior evaluation kernel for the Dirichlet problem consequently
has uniform angular smooth bounds and smooth center--radius dependence.
This follows directly by differentiating the Dirichlet equation and
using its uniform elliptic estimates; the center stays a fixed positive
distance from the unit sphere.  The chain rule for a base derivative
costs $|dr|=r\beta$, which is uniformly bounded by
\eqref{balls:bootstrap-condition}.

On the other cutoff set, $t>t_b$.  There the radii have a fixed positive
lower bound and all their required derivatives are bounded uniformly in
$\delta$.  The fixed-separation bounds of
Proposition~\ref{cross:initial}, followed by the smooth Dirichlet
construction of $\psi$, give the same conclusion.  Compactness of the
cutoff derivative supports completes the proof.
\end{proof}

Figure~\ref{fig:shrinking-functional} records the transfer and the two
parameter roles.  The next section establishes a bound for its density
that remains uniform as the spheres shrink.
\begin{figure}[htbp]
\centering
\begin{tikzpicture}[x=1cm,y=1cm,font=\small,>=Stealth,
  line cap=round,line join=round]
  \colorlet{sphereink}{blue!45!black}
  \node at (1.70,2.40) {One fixed spatial dilation $\lambda$};
  \draw[black!45,dashed] (1.70,0.55) circle (1.35);
  \draw[sphereink,thick] (1.70,0.55) circle (0.88);
  \draw[black!45,dashed] (1.70,0.55) circle (0.42);
  \fill (1.70,0.55) circle (1.6pt);
  \node[anchor=north east,inner sep=3pt] at (1.70,0.55) {$y$};
  \draw[->,sphereink] (1.70,0.55)--(2.3223,1.1723);
  \node[rotate=45,anchor=south,inner sep=3pt] at (2.011,0.861) {$r_\delta(y)$};
  \fill[sphereink] (2.58,0.55) circle (1.6pt);
  \draw[->,sphereink,thick] (2.66,0.55)--(4.02,0.55);
  \node[anchor=north,align=center,font=\footnotesize] at (3.42,0.38)
    {trace of\\$u_1$};
  \node[anchor=north,align=center] at (1.70,-1.02)
    {$t(y)<0$,\quad $\delta\downarrow0$\\[3pt]
     $r_\delta(y)\downarrow0$};
  \node[anchor=west,align=left,text width=7.10cm] at (4.30,0.55) {
    For a harmonic pair $(u_1,u_2)$:\\[7pt]
    $\displaystyle
      T_yu_1=\int_{S_y}\psi_\delta(y,\omega)\,
        u_1\!\left(\exp_y^\gamma(r_\delta(y)\omega)\right)
        \,d\sigma_y$\\[5pt]
    $\displaystyle \hphantom{T_yu_1}{}=u_2(y).$\\[9pt]
    Exact affine moments:\\[4pt]
    $\displaystyle T_y1=1,\qquad T_yx^i=y^i.$
  };
\end{tikzpicture}

\caption{A normal-coordinate section of $\gamma$-geodesic spheres in the
first variable, where $\gamma$ is the conformal reference metric.  The
center $y$ satisfies $t(y)<0$.  The spatial dilation $\lambda$ is fixed
while $\delta$ decreases.  The Green functional transfers paired harmonic
functions and reproduces affine harmonic coordinates exactly.  The
unit-sphere measure $d\sigma_y$ has total mass one; the density
$\psi_\delta$ need not be positive.}
\label{fig:shrinking-functional}
\end{figure}
